% ===== PRÉAMBULE CANONIQUE — à coller VERBATIM en tête de CHAQUE .tex =====
\documentclass[11pt,a4paper]{article}
\usepackage[utf8]{inputenc}\usepackage[T1]{fontenc}\usepackage{lmodern}
\usepackage[french]{babel}
\usepackage{amsmath,amssymb,mathtools}\usepackage{icomma}
\usepackage[version=4]{mhchem}          % \ce{...} : équations & formules
\usepackage{siunitx}\sisetup{locale=FR} % \qty{}{}, \si{}, \ang{}
\usepackage{chemfig}                    % \chemfig{...} : structures développées
\usepackage{xcolor}\usepackage{colortbl}
\usepackage{tikz}\usepackage{pgfplots}\pgfplotsset{compat=1.18}
\usepackage[most]{tcolorbox}\usepackage{enumitem}\usepackage{array,booktabs}
\usepackage[a4paper,margin=2.2cm]{geometry}
\usetikzlibrary{arrows.meta,calc,angles,quotes,positioning,patterns,backgrounds,shapes.geometric}
\definecolor{primary}{RGB}{222,49,99}\definecolor{primarydark}{RGB}{150,22,55}
\definecolor{defcol}{RGB}{0,114,178}\definecolor{methcol}{RGB}{52,92,148}
\definecolor{excol}{RGB}{0,135,90}\definecolor{attncol}{RGB}{200,40,55}
\tcbset{boxbase/.style={enhanced,breakable,boxrule=0.9pt,arc=2.5pt,
  left=3mm,right=3mm,top=2.3mm,bottom=2.3mm,fonttitle=\bfseries,coltitle=white}}
% --- boîtes TUTORIEL (noms EXACTS, ne pas renommer) ---
\newtcolorbox{cle}[1][]{boxbase,colback=defcol!6,colframe=defcol,
  title={\textsc{À retenir}\ifx\relax#1\relax\relax\else\textmd{ : \itshape #1}\fi}}
\newtcolorbox{methode}[1][]{boxbase,colback=methcol!7,colframe=methcol,sharp corners,
  title={\textsc{Méthode}\ifx\relax#1\relax\relax\else\textmd{ --- \itshape #1}\fi}}
\newtcolorbox{exemplebox}[1][]{boxbase,colback=excol!5,colframe=excol,
  title={\textsc{Exemple}\ifx\relax#1\relax\relax\else\textmd{ : \itshape #1}\fi}}
\newtcolorbox{piege}[1][]{boxbase,colback=attncol!6,colframe=attncol,
  title={\textsc{Piège}\ifx\relax#1\relax\relax\else\textmd{ : \itshape #1}\fi}}
\newtcolorbox{remarque}[1][]{boxbase,colback=black!4,colframe=black!45,
  title={\textsc{Remarque}\ifx\relax#1\relax\relax\else\textmd{ : \itshape #1}\fi}}
% --- boîtes EXERCICE / CORRIGÉ (noms EXACTS, auto-comptées) ---
\newtcolorbox[auto counter]{exo}[1][]{boxbase,colback=primary!4,colframe=primary,
  title={\textsc{Exercice}~\thetcbcounter\ifx\relax#1\relax\relax\else\textmd{ --- \itshape #1}\fi}}
\newtcolorbox[auto counter]{corrige}[1][]{boxbase,colback=excol!4,colframe=excol,
  title={\textsc{Corrigé}~\thetcbcounter\ifx\relax#1\relax\relax\else\textmd{ --- \itshape #1}\fi}}
\newcommand{\R}{\mathbb{R}}\newcommand{\N}{\mathbb{N}}\newcommand{\Z}{\mathbb{Z}}\newcommand{\ds}{\displaystyle}
\begin{document}
% THEME_TITLE: Nomenclature \& ordre de priorité des fonctions

Nommer une molécule, c'est d'abord choisir \emph{quelle} fonction lui donne son nom. Quand plusieurs fonctions coexistent, une seule est «\,principale\,» (elle donne le suffixe) ; les autres deviennent de simples préfixes.

\begin{cle}[L'échelle de priorité des fonctions]
De la plus prioritaire (à gauche) à la moins prioritaire :
\[
-\ce{COOH} > -\ce{COOR} > -\ce{CONH2} > -\ce{CHO} > -\ce{CO}- > -\ce{OH} > -\ce{NH2} > -\ce{O}- > \ce{C=C} > \ce{C#C}
\]
La fonction la plus haute présente est la \textbf{fonction principale} : c'est elle qui fixe le suffixe.
\end{cle}

\begin{cle}[Fonction principale (suffixe) vs fonction secondaire (préfixe)]
\begin{center}\small
\renewcommand{\arraystretch}{1.2}
\begin{tabular}{lll}
\toprule
fonction & suffixe (si principale) & préfixe (si secondaire) \\
\midrule
acide carboxylique & \textit{acide \dots-oïque} & carboxy- \\
aldéhyde & \textit{-al} & oxo- \\
cétone & \textit{-one} & oxo- \\
alcool & \textit{-ol} & hydroxy- \\
amine & \textit{-amine} & amino- \\
\bottomrule
\end{tabular}
\end{center}
Les \textbf{ramifications} et halogènes sont toujours des préfixes : méthyl-, éthyl-, chloro-, nitro-, méthoxy-\dots
\end{cle}

\begin{methode}[Nommer une molécule, pas à pas]
\begin{enumerate}[leftmargin=1.6em,itemsep=0.5pt]
  \item \textbf{Repérer} toutes les fonctions.
  \item \textbf{Classer} par priorité $\rightarrow$ la plus haute est principale (elle donne le suffixe) ; les autres passent en préfixes.
  \item \textbf{Choisir} la chaîne principale : la plus longue contenant la fonction principale.
  \item \textbf{Numéroter} pour donner le plus petit indice à la fonction principale.
  \item \textbf{Écrire} : préfixes (ordre alphabétique, avec indices) $+$ nom de la chaîne $+$ suffixe.
\end{enumerate}
\end{methode}

\begin{piege}[Le carbone fonctionnel porte le n$^\circ$1]
Pour l'\textbf{acide carboxylique} $-\ce{COOH}$, l'\textbf{ester} $-\ce{COOR}$, l'\textbf{aldéhyde} $-\ce{CHO}$ et l'\textbf{amide}, le carbone de la fonction porte \textbf{toujours} le numéro~1 : inutile d'écrire son indice. Pour les autres fonctions, on numérote simplement pour minimiser l'indice de la fonction principale.
\end{piege}

\end{document}
